\documentclass{beamer}
\usepackage[utf8]{inputenc}
\usepackage[main=spanish]{babel}

\usetheme{Madrid}
\useoutertheme[subsection=false]{miniframes}
\usecolortheme{default}

\uselanguage{spanish}
\languagepath{spanish}
\deftranslation[to=spanish]{Lemma}{Lema}
\deftranslation[to=spanish]{Theorem}{Teorema}

\setbeamertemplate{caption}[numbered]


\usepackage{nicefrac}
\usepackage{tikz}

% --- Configuración de enlaces ---

\let\originalhref\href
\renewcommand{\href}[2]{\originalhref{#1}{\color{blue}\ul{#2}}}

% --- Configuración de bibliografía ---
\usepackage[backend=biber]{biblatex} % 'style' puede ser numeric, alphabetic, etc.
\addbibresource{Bibliografia.bib} % Tu archivo de citas (.bib)



\newcommand{\bb}[1]{\mathbb{#1}}
\newcommand{\cc}[1]{\mathcal{#1}}
\newcommand{\ul}[1]{\underline{#1}}
\DeclareMathOperator{\diag}{diag}


\setbeamercolor{bloqueAmarilloCc}{bg=orange!90!yellow, fg=white} % Cabecera
\setbeamercolor{bloqueAmarilloCu}{bg=orange!6!yellow!4!white, fg=black} % Cuerpo

\newenvironment{definitionblock}[1]{%
  \begin{beamerboxesrounded}[upper=bloqueAmarilloCc, lower=bloqueAmarilloCu, shadow=true]{#1}%
}{%
  \end{beamerboxesrounded}%
}


\AtBeginSection[]
{
  \begin{frame}
    \frametitle{Índice}
    \tableofcontents[currentsection]
  \end{frame}
}


% Declarar aplicaciones
% 1. Nombre aplicación
% 2. Dominio
% 3. Codominio
% 4. Variable
% 5. Imagen de la variable
\newcommand{\Func}[5]{
    \begin{equation*}
        \begin{array}{rrll}
            \displaystyle #1:& \displaystyle  #2 & \longrightarrow & \displaystyle  #3\\
               & \displaystyle  #4 & \longmapsto & \displaystyle  #5
        \end{array}
    \end{equation*}
}
\DeclareMathOperator{\tr}{Tr}


\usepackage[makeroom]{cancel}






\title[3.3. DSM]{Capítulo 3.3:\\Denoising Score Matching}
\subtitle{Trabajo de Fin de Grado}
\author[Arturo]{Arturo Olivares Martos}
\institute[DGIIM - UGR]{Doble Grado en Ingeniería Informática y Matemáticas\\Universidad de Granada}
\date{\today}
\logo{\includegraphics[height=1cm]{Aux/UGR-Logo.png}}




\begin{document}

\begin{frame}
    \titlepage
\end{frame}

\section{Introducción}
\begin{frame}{Motivación}
    
    El Sliced Score Matching (SSM) está bien, pero...
    \begin{itemize}
        \item Normalmente, $p_{\text{data}}$ es nula en la mayoría de los puntos del espacio. Al calcular $\nabla_{x}\log p_{\text{data}}(x)$, nos encontramos con que la mayoría de los puntos no tienen información útil: no está definido, o es inestable.\pause
        
        En el muestreo de Langevin, ¡¡no sabemos hacia dónde movernos!!\pause
        \item JVP y VJP, aunque evitan calcular $D$ derivadas parciales, siguen siendo costosos computacionalmente e inducen errores de aproximación.
    \end{itemize}
\end{frame}

\begin{frame}{Solución: Inyección de ruido}
    En vez de aproximar $s_{\text{data}}(x)$, aproximamos $s_{\tilde{p}}(x)$, donde $\tilde{p}(x)$ es una versión modificada de $p_{\text{data}}(x)$. \pause
    \begin{itemize}
        \item Fijado un nivel de ruido $\sigma$, inyectamos ruido a los datos empleando $p_{\sigma}(\tilde{x}|x)$.\pause
        \item El modelo $s_{\phi}(\wt{x},\sigma)$ aproxima $s_{p_\sigma}(\tilde{x})$, el score de la marginal de los datos ruidosos:
        \begin{equation*}
            p_\sigma(\tilde{x}) = \int_{\mathbb{R}^D} p_{\sigma}(\tilde{x}|x)p_{\text{data}}(x)dx
        \end{equation*}\pause
        \item Si $\sigma\approx 0$, $p_\sigma(\tilde{x})\approx p_{\text{data}}(x)$ y $s_{p_\sigma}(\tilde{x})\approx s_{\text{data}}(x)$. Por tanto, podremos emplearlo para generar muestras de $p_{\text{data}}(x)$ mediante muestreo de Langevin.
    \end{itemize}
\end{frame}

\begin{frame}{Ventajas de la inyección de ruido}
    $p_\sigma(\tilde{x})$ es una versión suavizada de $p_{\text{data}}(x)$. Si $p_{\sigma}(\tilde{x}|x)=\cc{N}(\tilde{x};x,\sigma^2\mathbf{I})$, entonces:
    \begin{equation*}
        p_\sigma = p_{\text{data}} * \cc{N}(0,\sigma^2\mathbf{I})
    \end{equation*}\pause
    \begin{itemize}
        \item $p_\sigma(\tilde{x})$ es más estable y está definido en todo el espacio, por lo que $s_{p_\sigma}(\tilde{x})$ también lo está.\pause
        \item Mayor cobertura del espacio de datos, pues es más sencillo para el muestreo de Langevin moverse entre regiones de alta densidad de $p_\sigma(\tilde{x})$.
    \end{itemize}    
\end{frame}

\begin{frame}{Efectos de la elección de $\sigma$}
    Trade-off entre estabilidad y precisión:\pause
    \begin{itemize}
        \item Si $\sigma$ es demasiado grande, $p_\sigma(\tilde{x})$ es demasiado suave y no se parece a $p_{\text{data}}(x)$.
        
        Es útil para explorar el espacio de datos, pero tan solo aporta estructuras globales de $p_{\text{data}}(x)$.\pause
        \item Si $\sigma$ es demasiado pequeño, $p_\sigma(\tilde{x})$ se parece a $p_{\text{data}}(x)$, pero es inestable y no está definido en todo el espacio.
        
        Proporciona información más detallada sobre la estructura de los datos, pero puede ser difícil de usar en la práctica.
    \end{itemize}\pause

    \begin{center}
        ¿Necesitamos un único nivel de ruido $\sigma$? \pause {\color{green!50!black} ¡No!}\pause\\
        Pero esos serán los NCSN, tranquilidad...
    \end{center}
\end{frame}


\section{Entrenamiento}

\begin{frame}{Entrenamiento}
    \begin{equation*}
        \cc{L}_{\text{SM}}(\phi,\sigma) := \frac{1}{2}\bb{E}_{\tilde{x}\sim p_\sigma(\tilde{x})}[\|s_\phi(\tilde{x},\sigma) - s_{p_\sigma}(\tilde{x})\|_2^2]
    \end{equation*}\pause\vspace{-0.8cm}
    \begin{center}
        {\color{red} ¡Intractable!}
    \end{center}
    \pause
    \begin{theorem}[Equivalencia entre $\cc{L}_\text{DSM}$ y $\cc{L}_\text{SM}$]
        Fijado un nivel de ruido $\sigma$, se cumple que:
        \begin{align*}
            \cc{L}_\text{SM}(\phi,\sigma) &= \cc{L}_\text{DSM}(\phi,\sigma) + C_\sigma\\
            \cc{L}_\text{DSM}(\phi,\sigma) :&= \frac{1}{2}\bb{E}_{x\sim p_{\text{data}}(x), \tilde{x}\sim p_\sigma(\tilde{x}|x)}[\|s_\phi(\tilde{x},\sigma) - s_{p_\sigma}(\tilde{x}|x)\|_2^2]
        \end{align*}\pause
        El valor óptimo es:
        \begin{equation*}
            s_\phi^*(\tilde{x},\sigma) = s_{p_\sigma}(\tilde{x}) = \nabla_{\tilde{x}}\log p_\sigma(\tilde{x})
        \end{equation*}
    \end{theorem}
\end{frame}

\begin{frame}{Ejemplo: $p_\sigma(\tilde{x}|x) = \cc{N}(\tilde{x};x,\sigma^2\mathbf{I})$}

    Un caso muy común es el de inyectar ruido gaussiano isotrópico a los datos:
    \begin{equation*}
        \wt{x} = x + \sigma \epsilon, \quad \epsilon \sim \cc{N}(0,\mathbf{I})
    \end{equation*}\pause
    \begin{equation*}
        s_{p_\sigma}(\tilde{x}|x) = \nabla_{\tilde{x}}\log p_\sigma(\tilde{x}|x) = \frac{x-\tilde{x}}{\sigma^2}
    \end{equation*} \pause
    \begin{align*}
        \cc{L}_\text{DSM}(\phi,\sigma) &= \frac{1}{2}\bb{E}_{x\sim p_{\text{data}}(x), \tilde{x}\sim p_\sigma(\tilde{x}|x)}[\|s_\phi(\tilde{x},\sigma) - s_{p_\sigma}(\tilde{x}|x)\|_2^2] \\
        &= \frac{1}{2}\bb{E}_{x\sim p_{\text{data}}(x), \epsilon\sim \cc{N}(0,\mathbf{I})}\left[\|s_\phi(x+\sigma\epsilon,\sigma) + \frac{\epsilon}{\sigma}\|_2^2\right]
    \end{align*}\pause
    \begin{center}
        {\color{green!50!black} ¡Tractable!}
    \end{center}
\end{frame}

\section{Por qué ``Denoising''}
\begin{frame}{Por qué ``Denoising''}

    Los modelos estudiados se denominan ``Denoising'' Score Matching, pero hasta ahora tan solo hemos hablado de inyectar ruido a los datos.\pause
    \begin{lemma}[Fórmula de Tweedie]
        Sea $p_\sigma(\tilde{x}|x) = \cc{N}(\tilde{x};\alpha x,\sigma^2\mathbf{I})$, con $\alpha\footnote{En nuestro caso, $\alpha=1$.}\in\mathbb{R}^*$. Entonces:
        \begin{equation*}
            \alpha \bb{E}[x|\tilde{x}] = \tilde{x} + \sigma^2 s_{p_\sigma}(\tilde{x}) = \tilde{x} + \sigma^2 \nabla_{\tilde{x}}\log p_\sigma(\tilde{x})
        \end{equation*}
    \end{lemma}
    
    $E[x|\tilde{x}]$ es el estimador de mínima varianza de $x$ dado $\tilde{x}$, por lo que se puede entender como una forma de ``denoising'' de $\tilde{x}$.\pause

    \begin{center}
        ¿Podemos emplearlo de forma directa como \emph{denoiser}?\pause\\
        No. Como en los DDPM, puede ser una distribución multimodal, y el estimador $\bb{E}[x|\tilde{x}]$ puede no pertenecer a ninguna de las modas.
    \end{center}


\end{frame}

    






\begin{frame}[allowframebreaks]{Bibliografía}
    \nocite{lai2026principlesdiffusionmodels}
    \printbibliography
\end{frame}




\end{document}